% ===================================================================
% SECTION 3: METHODOLOGY
% File: sections/methods003.tex
% ===================================================================

\subsection{Parametric Time-Series Deformation Modeling}
\label{subsec:parametric_deformation}

Multi-source observation series recorded across the Choushui River Alluvial Fan exhibited complex superpositions of long-term subsidence trends, seasonal oscillations, and sudden instrument offsets. Decomposing raw observation series required a uniform parametric signal decomposition model applicable across subsurface extensometer records, continuous GNSS (cGNSS) position series, and SBAS-InSAR line-of-sight displacement series \citep{hung_exploring_2024}. The observational time series $y(t)$ at a given monitoring point was modeled as a linear combination of baseline offset, constant velocity, seasonal harmonics, and discrete step jumps \citep{nikolaidis_observation_2002}
\begin{equation}
y(t) = a + b \, t + \sum_{m=1}^{M} \left[ d_m \sin\left(2\pi f_m t\right) + e_m \cos\left(2\pi f_m t\right) \right] + \sum_{j=1}^{J} o_j H(t - T_j) + \varepsilon(t)
\label{eq:parametric_deformation}
\end{equation}
where $a$ represents the initial datum offset, $b$ denotes the linear displacement velocity, $M$ specifies the number of seasonal harmonic frequencies, $d_m$ and $e_m$ represent sine and cosine harmonic amplitudes, $J$ indicates the total number of discrete step offsets, $o_j$ specifies the magnitude of the $j$-th offset occurring at time $T_j$ modeled via the Heaviside step function $H(t - T_j)$, and $\varepsilon(t)$ represents unmodeled observation noise.

Seasonal periodic oscillations were parameterized using annual ($f_1 = 1~\text{yr}^{-1}$) and semiannual ($f_2 = 2~\text{yr}^{-1}$) frequencies to isolate annual groundwater pumping drawdowns and monsoon recharge cycles. Equipment replacements, sensor recalibrations, and physical wellhead maintenance events were isolated through step function terms $o_j H(t - T_j)$ at known event epochs $T_j$. Least-squares estimation derived optimal model parameters ($a, b, d_m, e_m, o_j$) for each observation point. Removal of step offset jumps $o_j$ and seasonal terms constructed cleaned displacement series, enabling monthly first-differencing within six uniform 50~m depth sections (designated as S1 through S6, spanning 0--300~m depth) to derive depth-dependent compaction rates down the well profile.

\subsection{Compositional Sediment Transformation via Isometric Log-Ratio and Sequential Binary Partition}
\label{subsec:ilr_sbp_transformation}

Sedimentological compositions derived from borehole lithological logs and 3D hydrogeological modeling established the static skeletal storage framework controlling physical compressibility across the six 50~m depth sections. High-resolution lithological profiles recorded volumetric percentages of four sedimentary facies, comprising gravel, coarse sand, fine sand, and mud (silt and clay). Because raw volumetric fractions sum to 100 percent within any given depth section, direct inclusion of raw percentages induced constant-sum closure and severe multicollinearity in linear statistical estimators.

A Sequential Binary Partition (SBP) hierarchy was constructed to map the four sediment fractions into three unconstrained, orthogonal log-ratio balances. The SBP hierarchy partitioned the sediment components sequentially: the first balance ($ILR_1$) contrasted coarse-grained facies (gravel and coarse sand) against fine-grained facies (fine sand and mud); the second balance ($ILR_2$) contrasted gravel against coarse sand within the coarse sub-group; the third balance ($ILR_3$) contrasted fine sand against mud within the fine sub-group. The Isometric Log-Ratio (ILR) transformation for a binary partition splitting $r$ components in group 1 and $s$ components in group 2 was calculated as \citep{egozcue_isometric_2003, oh_using_2024}
\begin{equation}
ILR_k = \sqrt{\frac{r \cdot s}{r + s}} \ln \left( \frac{\left( \prod_{j \in \text{Group 1}} x_j \right)^{1/r}}{\left( \prod_{m \in \text{Group 2}} x_m \right)^{1/s}} \right)
\label{eq:ilr_transformation}
\end{equation}
This ILR transformation mapped the four constrained volume fractions into three orthogonal coordinates ($ILR_{1}, ILR_{2}, ILR_{3}$) for each 50~m depth section, eliminating constant-sum closure and providing unconstrained sedimentological predictors for layerwise compaction modeling.

\subsection{Formulation of Hydrogeological Drivers}
\label{subsec:hydrogeological_drivers}

Subsurface aquifer-system compaction is physically driven by effective stress changes, which are directly proportional to groundwater level drawdowns. However, piezometer wells across the monitoring network are screened at discrete aquifer intervals, whereas MLCW measurements were aggregated into six uniform 50~m depth sections (designated as S1 through S6, spanning 0--300~m depth) to establish consistent depth intervals across boreholes. Assigning piezometric head variations to these 50~m sections required resolving vertical screen depth mismatches. For depth sections containing co-located piezometer wells whose screen midpoint fell within the section depth interval expanded by 10~m, direct monthly piezometric head series were assigned to drive that section. For unmonitored depth sections lacking qualified co-located wells, head series were spatially interpolated at the exact extensometer station coordinates from surrounding regional piezometers using Ordinary Kriging applied to screen-depth-matched regional wells. Monthly first-differencing ($\Delta H_{i,s}(t) = H_{i,s}(t) - H_{i,s}(t-1)$) isolated transient monthly drawdowns driving effective stress rates $\Delta \sigma'(t) = -\gamma_w \Delta H(t)$ while removing arbitrary wellhead elevation offsets \citep{galloway_land_1999}.

Subsurface compaction depends on both extrinsic dynamic forcing and intrinsic skeletal storage properties \citep{liu_reconstructing_2023}. Low vertical hydraulic conductivity within fine-grained silt and clay aquitards restricts pore-water dissipation rates, causing compaction to lag months or years behind hydraulic head drawdowns in adjacent aquifers. To capture this delayed consolidation memory, discrete monthly head-change lags ($\Delta H_{i,s}(t - \tau)$ for $\tau \in \{1, 3, 6, 12, 24\}$~months) were constructed across a multi-year horizon. Seasonal agricultural pumping intensity was incorporated by modulating hydraulic lags during dry-season months ($k \times \text{DrySeason}$). Intrinsic skeletal compressibility constraints were coupled by interacting hydraulic lags with clay content balances ($\Delta H \times ILR_{\text{clay}}$). Top-boundary mechanical constraints above the multi-layer system were provided by surface acceleration $d^2S_{\text{total}}(t)$, computed as the second-difference of integrated surface displacement series derived from combined cGNSS and SBAS-InSAR observations.

\subsection{Bayesian Ridge Regression Estimation Engine}
\label{subsec:bayesian_ridge_engine}

The estimation engine was constructed using Bayesian Ridge Regression (BRR) to resolve severe multicollinearity among lagged hydraulic head predictors while preventing over-fitting. High-dimensional predictor matrices containing 24-month discrete lags and cross-section interaction terms exhibit strong collinearity, rendering standard Ordinary Least Squares (OLS) estimators highly unstable. The BRR framework applies a zero-mean isotropic Gaussian prior over model coefficient vectors $\boldsymbol{w}_s \in \mathbb{R}^P$ for each section $s$
\begin{equation}
p(\boldsymbol{w}_s \mid \alpha) = \mathcal{N}\left( \boldsymbol{w}_s \mid \mathbf{0}, \alpha^{-1}\mathbf{I}_P \right)
\label{eq:weight_prior}
\end{equation}
where $\alpha$ represents the precision hyperparameter of the weight prior and $\mathbf{I}_P$ denotes the $P \times P$ identity matrix.

Model likelihood was formulated assuming Gaussian observation noise with precision parameter $\lambda$
\begin{equation}
p(\mathbf{y}_s \mid \mathbf{X}_s, \boldsymbol{w}_s, \lambda) = \mathcal{N}\left( \mathbf{y}_s \mid \mathbf{X}_s \boldsymbol{w}_s, \lambda^{-1}\mathbf{I}_N \right)
\label{eq:likelihood}
\end{equation}
Combining likelihood and prior via Bayes' theorem derived an analytical Gaussian posterior distribution over weight vectors $p(\boldsymbol{w}_s \mid \mathbf{X}_s, \mathbf{y}_s, \alpha, \lambda) = \mathcal{N}\left( \boldsymbol{w}_s \mid \boldsymbol{\mu}_s, \boldsymbol{\Sigma}_s \right)$, where posterior covariance $\boldsymbol{\Sigma}_s$ and mean $\boldsymbol{\mu}_s$ were calculated as
\begin{equation}
\boldsymbol{\Sigma}_s = \left( \lambda \mathbf{X}_s^T \mathbf{X}_s + \alpha \mathbf{I}_P \right)^{-1}, \quad \boldsymbol{\mu}_s = \lambda \boldsymbol{\Sigma}_s \mathbf{X}_s^T \mathbf{y}_s
\label{eq:posterior_distribution}
\end{equation}
Hyperparameters $\alpha$ and $\lambda$ were jointly estimated from observational data via empirical Bayes marginal likelihood maximization using L-BFGS optimization, providing automatic $L_2$ weight regularization and analytical predictive variance estimates without manual hyperparameter tuning.

\subsection{Multi-Tier Operational Validation Framework}
\label{subsec:operational_validation_framework}

Model performance and generalizability were evaluated across three complementary validation tiers designed to test temporal stability, spatial transferability, and operational delay resilience. Temporal stability was assessed using walk-forward cross-validation, where models were trained on historical records up to month $t$ and evaluated on subsequent held-out test windows. Evaluation metrics comprised the Coefficient of Determination ($R^2$), Root Mean Square Error ($\text{RMSE}$), and Mean Absolute Error ($\text{MAE}$) computed separately for each 50~m depth section.

Spatial transferability across unmonitored fan locations was evaluated using Leave-One-Station-Out (LOSO) cross-validation. In each LOSO iteration, all data from one primary MLCW station (such as Tuku) were completely withheld from model training, and the model fit on the remaining four stations was evaluated on the held-out station. Pooled and section-specific $R^2$ metrics on held-out stations validated whether learned hydrogeological weight parameters transferred effectively across the alluvial fan system without station-specific recalibration.

Uncertainty quantification was implemented via Split-Conformal Prediction to generate non-parametric 90\% prediction intervals with distribution-free coverage guarantees. Calibration residuals from a separate validation split were used to compute empirical quantiles $\hat{q}_{0.90}$, deriving lower and upper estimation bounds $\hat{C}(\mathbf{x}_t) = [\hat{y}_t - \hat{q}_{0.90}, \hat{y}_t + \hat{q}_{0.90}]$ that encapsulated predictive uncertainty without assuming Gaussian residual errors.

Operational resilience under realistic inspection delays was evaluated using a 3-arm backtest framework simulating a 6-month label update delay. The backtest benchmarked three operational modes: (1) a \textit{Refit Arm} updated immediately as new monthly extensometer readings became available, (2) a \textit{Frozen Arm} trained on historical baseline data and operating without weight updates during the 6-month label delay window, and (3) a \textit{Naive Persistence Arm} assuming monthly compaction equaled the previous month's value. Performance comparisons across the three arms quantified model staleness decay and confirmed estimation reliability under field monitoring constraints.